\documentclass[a4paper, 11pt]{article}

% -- Coments
\usepackage{verbatim}

% ----- Fonts -----
% -- Fuente --
% \usepackage{fontspec}
% \setmonofont{JetBrainsMono Nerd Font}

\usepackage{babel}
\babelprovide[import, main]{spanish}

% -- Color --
\usepackage{xcolor}
%\definecolor{azul}{RGB}{00,33,99}
\definecolor{azul}{RGB}{35,72,180}

% -- Page Margin --
\usepackage[margin=1in, headheight=15pt]{geometry}

% -- Encabezado con nombre y expediente --
\usepackage{fancyhdr}
\pagestyle{fancy}
\fancyhf{}
\lhead{\textbf{Bórquez Guerrero Angel Fernando}}
\rhead{Expediente: 219208106}
\renewcommand{\headrulewidth}{0.5pt}

% -- Espaciados --
\newcommand{\Pspace}{0.5cm}
\newcommand{\Aspace}{0.2cm}

% -- Columnas --
\usepackage{multicol}

% -- Imagenes --
\usepackage{graphicx}
\usepackage{float}

% -- Matemáticas --
\usepackage{amsmath, amssymb}

% -- Gráficas --
\usepackage{pgfplots}
\pgfplotsset{compat=1.18}


\title
{
    \textbf{Contexto Regional 2026-2} \\[0.3cm]
    {\large \textit{Los desastres y su influencia en el ámbito político: el incendio de la Guardería ABC en las elecciones en el estado de Sonora, 2009}} \\[0.2cm]
    {\normalsize Carrillo Luvianos, Mario Alejandro; Toscana Aparicio, Alejandra y Martínez Rangel, Rubí} \\[0.3cm]
    {\normalsize Profesor José Moreno Vega}
}
\date{}

\begin{document}

\begin{otherlanguage}{spanish}
\maketitle
\end{otherlanguage}
\thispagestyle{fancy}
\vspace{-0.5cm}


{\Large 1. ¿Cuál es el argumento central del autor?} \\
{\color{azul}
    Los autores dicen que el incendio de la Guardería ABC, donde murieron 49 niños el 5 de junio de 2009 en Hermosillo, afectó mucho al PRI en las elecciones de ese año y aceleró el cambio político en Sonora. Explican que un desastre se vuelve político cuando la gente siente que se pudo evitar y que el gobierno lo manejó mal, y eso fue lo que pasó: los gobiernos federal, estatal y municipal se culparon entre ellos, y la indignación aumentó cuando se supo que los dueños de la guardería eran cercanos al gobernador. Por eso el PAN le ganó al PRI la gubernatura y la presidencia municipal de Hermosillo. Eso sí, aclaran que el incendio no fue la única causa, porque el PRI ya estaba dividido por la imposición de su candidato y la competencia venía creciendo desde 1997, aunque el efecto fue mayor en el municipio y en el congreso local que en la gubernatura.
}
\vspace{1cm}

{\Large 2. ¿Qué evidencias utiliza para respaldarlo?} \par
{\color{azul}
Para empezar compara con otros desastres que terminaron cambiando la política, como el ciclón de Bangladesh, el sismo de Nicaragua, Chernobyl, el huracán Katrina y, en México, los sismos de 1985, las explosiones de Guadalajara y el huracán Paulina en Acapulco. Después cuenta lo que pasó con el incendio casi día por día, desde la demanda del IMSS y las acusaciones entre los gobiernos, hasta la marcha de 20 mil personas del 4 de julio y la orden de aprehensión contra los dueños el mismo día de la elección. También usa los números del Consejo Estatal Electoral, con un cuadro que compara los votos de los tres partidos en las elecciones para gobernador de 1991 a 2009, y se apoya en autores como Olson y Gawronski, López-Levi y Poom.
}
\vspace{1cm}

{\Large 3. ¿Cómo se relaciona esta lectura con la lectura anterior?} \par
{\color{azul}
Las dos hablan de la cercanía entre el gobierno de Sonora y las empresas privadas, y de cómo esto afectó a la gente común. En la lectura 11 el gobierno daba exenciones de impuestos y dejó el parque industrial en manos de inversionistas, y aquí el IMSS entregó el manejo de las guarderías a particulares cercanos al gobernador. En ambos casos las consecuencias las vivieron los trabajadores con la crisis de 1974 y los niños y sus familias con el incendio.
}

\end{document}
